\section{The Question and the Method}\label{sec:q}

The question under test: \emph{when a team wins the Presidential Hackathon's
International Track, is it because the project is genuinely excellent --- or
because it lands on the nation's immediate need at the right moment?} The
question matters because our own entry strategy depends on the answer.

\textbf{Method.} We identified the publicly named Teams of Excellence for each
International Track edition, 2019--2025 (two per year in official reporting;
14 in total), then researched each against three factual axes: (1) what the
project actually was and how technically substantial it was \emph{at the time
of winning}; (2) whether it pre-existed the competition; (3) whether any
public trace of activity exists \emph{after} the competition (websites, code
repositories, partnerships, press, return appearances). Sources are official
press releases (MODA, Office of the President), the competition's own history
pages, the winning organizations' own publications, and third-party records
(Open Contracting Partnership, CoST International, GitHub). Where we found
nothing, we say ``no public trace found'' --- which is evidence of absence
only in the weak sense, and is labeled accordingly.

\textbf{Two bias controls.} First, this review deliberately re-verified claims
made in this project's own earlier research files, and corrects two of them
(Section~\ref{sec:corrections}). Second, we observe winners only: the losing
pool is invisible, so we cannot prove that better projects lost. Claims below
are therefore about \emph{what winners have in common}, not about what judges
rejected.

\textbf{The scoring frame the judges use.} The rubric weights Feasibility 40\%,
Innovation 30\%, Social Impact 30\% at preliminaries; finals add
Implementation \& Verification at 30\%\cite{handbook}. Note what this
arithmetic already implies: ``uniqueness of technology'' is at most
three-tenths of the decision; the other seven-tenths reward a solvable,
important, demonstrable problem. The theme itself is set annually to the
presidential agenda --- 2026's theme ``echoes President Lai's New Year
vision''\cite{handbook} --- and the award is presented personally by the
President in a ceremony attended by diplomatic representatives\cite{focuscrop}.
The institutional design, in other words, is not neutral between our two
hypotheses. What follows tests whether the record matches the design.

\section{The Case Files, 2019--2025}\label{sec:cases}

Table~\ref{tab:master} compresses the fourteen case files; the paragraphs
after it add what a table cannot.

\begin{table}[htbp]
\caption{Fourteen Teams of Excellence: what they were, whether they pre-existed, what remains}
\label{tab:master}
\footnotesize
\begin{tabularx}{\textwidth}{@{}p{0.032\textwidth}p{0.145\textwidth}p{0.105\textwidth}Xp{0.205\textwidth}@{}}
\toprule
\textbf{Yr} & \textbf{Team} & \textbf{Origin} & \textbf{Project (verified description)} & \textbf{Pre-existence / afterlife} \\
\midrule
19 & CoST Honduras & Honduras & SISOCS infrastructure-disclosure platform; disclosure rates 27\%$\to$95\%\cite{cost} & Pre-existed (2014--15, World Bank context); still operating, 1,850+ projects, $\sim$US\$1bn disclosed\cite{cost} \\
19 & Cartelogy & Malaysia & Procurement red-flag tool; $\sim$7,000 contracts converted to OCDS\cite{ocpcartel} & Backed by civic-tech group; no public post-2019 trace found \\
20 & Fiscal Force & India & Health Procurement Performance Index on Himachal Pradesh tenders\cite{ocpfiscal} & Built by CivicDataLab (existing org); grew into a formal CivicDataLab--OCP partnership with state governments\cite{ocpindia} \\
20 & Learning Man & Taiwan & Social-housing procurement mapping via OC4IDS\cite{ocpfiscal} & No public post-2020 trace found \\
21 & A/B Street & Intl.\ (US) & Open-source traffic simulation / street redesign tool\cite{abstreet} & Pre-existed (2018); post-win: creator ``almost ended the project,'' continued at the Alan Turing Institute, now a one-person consultancy; development narrowed to a library\cite{abstreet,abinterview} \\
21 & Office Farmers & Intl. & Open-data crop adaptation modeling\cite{pres2021} & No public post-2021 trace found \\
22 & BRSDM@Er-Lin & Taiwan & Field biochar-pyrolysis equipment + cloud data platform, promoted toward Southeast Asia\cite{hist2023} & Research-team project; no public post-2022 trace found \\
22 & B.E.N.Z. & India & GIS-based urban-forest and tree-placement optimization\cite{hist2023} & No public post-2022 trace found \\
23 & HysonTech & Taiwan/India & ``FarmABetterFish'' AIoT aquaculture system\cite{moda2023} & Hyson Technology sells intelligent-aquaculture products commercially today\cite{hyson} \\
23 & Re-Fill City & Thailand & App incentivizing bottle refill/reuse with retailer partnerships\cite{moda2023} & No public post-2023 trace found \\
24 & GreenhopeBCTW & Taiwan/US & Blockchain personal carbon wallet; monetizes verified behavior\cite{moda2024} & GreenHope operates as a Web3 company today\cite{greenhope} \\
24 & MooApps & Indonesia & Digital livestock-monitoring app\cite{moda2024} & No public post-2024 trace found \\
25 & CropNow & India & NEST IoT device + CropDesk advisory: soil, microclimate, disease prediction\cite{focuscrop} & Invited back to the 2026 launch to present implementation progress\cite{moda2026} \\
25 & Beyond Hearing & Taiwan/France & AI+AR real-time sound-to-visual translation for the hearing-impaired\cite{pres2025} & Invited back to the 2026 launch to present implementation progress\cite{moda2026} \\
\bottomrule
\end{tabularx}
\end{table}

\subsection{What the winners were, in their year's context}

\textbf{2019--2020: the transparency years.} The theme was sustainable
infrastructure; the global civic-tech zeitgeist was open contracting. Both
years' winners were procurement-data tools. Neither was algorithmically novel
--- SISOCS is a disclosure database; HPPI is an index over tender data --- but
both were \emph{real}: SISOCS had been disclosing actual road projects since
2015 under a World Bank program\cite{cost}, and Fiscal Force was built by an
established data-for-governance organization\cite{ocpindia}. The technically
interesting entrant of 2019, Cartelogy's red-flag cartel detection, is the one
whose trail goes cold afterward.

\textbf{2021--2022: the net-zero years.} Themes pivoted to climate action.
A/B Street won: a genuinely impressive piece of software --- but one started
in 2018, three years mature, with an existing community\cite{abstreet}. Its
post-win account is the most instructive in the whole record: the creator
nearly shut it down within a year of winning, and the project survived by
institutional adoption (Turing Institute) unrelated to
Taiwan\cite{abinterview}. The 2022 winners (field biochar pyrolysis; GIS tree
placement) were competent applied projects squarely on the ``practicing net
zero'' theme\cite{hist2023}; neither left a public trace we could find.

\textbf{2023--2025: the AI-for-the-vulnerable years.} As generative AI
exploded, winning entries became AIoT systems with named beneficiary
populations: fish farmers (HysonTech --- an operating company whose product
line exists today\cite{hyson}), consumers (Re-Fill City), individual citizens'
carbon behavior (GreenhopeBCTW --- likewise an operating
company\cite{greenhope}), smallholder farmers (CropNow), and the
hearing-impaired (Beyond Hearing). The 2025 pair --- an Indian agri-IoT
startup-style team and a Taiwan--France accessibility team --- were selected
from nine finalists and honored by the President alongside India's
representative office chief\cite{focuscrop}.
